\documentclass[11pt]{article}
\usepackage{amsmath,amssymb,amsthm,booktabs,hyperref}
\newtheorem{theorem}{Theorem}
\newtheorem{lemma}{Lemma}
\newtheorem{proposition}{Proposition}
\newtheorem{definition}{Definition}

\title{Random Access to Fixed-Prime Smooth Numbers via Exponent-Stream Self-Similarity}
\author{Mihai Polceanu and collaborators}
\date{Draft skeleton}

\begin{document}
\maketitle

\begin{abstract}
We study the fixed-prime smooth-number random-access problem: given a finite set of primes $P=\{p_1,\ldots,p_k\}$ and a rank $N$, compute the exponent vector of the $N$-th integer of the form $p_1^{e_1}\cdots p_k^{e_k}$. We describe a lattice formulation, an exponent-stream self-similarity principle, and a zero-position interpretation in terms of lower-dimensional faces. These observations yield a fast three-prime ranker based on two-prime floor sums and a layer-compressed meet-in-the-middle method for five primes. We present certified rank checks for selected five- and six-prime instances and compare against the standard dynamic-programming/pointer algorithm for super ugly numbers and practical meet-in-the-middle baselines.
\end{abstract}

\section{Introduction}
Define the problem, distinguish fixed-prime random access from sequential generation and from general $y$-smooth counting.

\section{Background and related work}
Discuss Hamming/regular numbers, super ugly number DP algorithms, smooth-number counting, $X+Y$ selection, and fixed-dimensional lattice counting. The working related-work notes are maintained in \texttt{paper/related\_work.md}; in particular, separate implemented comparators from unimplemented literature baselines such as Frederickson--Johnson sorted-matrix selection, soft-heap $X+Y$ selection, Barvinok-style fixed-dimensional lattice counting, and Bernstein-style smooth-part algorithms.

\section{The log-lattice model}
\begin{definition}
For a finite prime set $P$, let $S_P$ be the ordered sequence of $P$-smooth integers, with $1$ first.
\end{definition}

\begin{lemma}[Log-lattice equivalence]
Ordering $P$-smooth numbers is equivalent to ordering exponent vectors $e\in\mathbb N^k$ by $\sum_i e_i\log p_i$.
\end{lemma}

\section{Exponent streams and face recursion}
\begin{definition}
Let $E_i(n)$ be the exponent of $p_i$ in the $n$-th $P$-smooth number.
\end{definition}

\begin{theorem}[Exponent-stream self-similarity]
For each $i$, deleting zeroes from $E_i$ and subtracting $1$ from each remaining term yields $E_i$ again.
\end{theorem}
\begin{proof}
The positive part of the $p_i$ exponent stream corresponds exactly to numbers divisible by $p_i$. Dividing each such number by $p_i$ maps them bijectively and order-preservingly onto $S_P$.
\end{proof}

\begin{theorem}[Zeroes as ranked faces]
The zero positions in $E_i$ are the ranks in $S_P$ of the $P\setminus\{p_i\}$-smooth numbers.
\end{theorem}

\section{Two- and three-prime rank counting}
State the floor-sum formula for $C_{p,q}(T)$ and the resulting $C_{p,q,r}(T)$ algorithm.

\section{Layer-compressed meet-in-the-middle}
Give the five-prime decomposition and quotient/residue count identity.

\section{Certification}
Describe interval-log rank auditing and exact ambiguity resolution.

\section{Benchmarks}
Include certified tables and comparisons.

\section{Limitations}
Not a general $\Psi(x,y)$ algorithm; not a prime-generation method; fastest high-$k$ variants require further certification.

\bibliographystyle{plain}
\begin{thebibliography}{9}
\bibitem{bernstein-smoothparts} D. J. Bernstein, How to find smooth parts of integers, 2004.
\bibitem{johnson-mizoguchi} D. B. Johnson and T. Mizoguchi, Selecting the Kth Element in $X+Y$ and $X_1+X_2+\cdots+X_m$, SIAM Journal on Computing, 7(2), 1978.
\bibitem{mirzaian-arjomandi} A. Mirzaian and E. Arjomandi, Selection in $X+Y$ and matrices with sorted rows and columns, Information Processing Letters, 20(1), 1985.
\bibitem{fj-xplusy} G. N. Frederickson and D. B. Johnson, The Complexity of Selection and Ranking in $X+Y$ and Matrices with Sorted Columns, Journal of Computer and System Sciences, 24(2), 1982.
\bibitem{fj-sorted-matrices} G. N. Frederickson and D. B. Johnson, Generalized Selection and Ranking: Sorted Matrices, SIAM Journal on Computing, 13(1), 1984.
\bibitem{kaplan-softheap} H. Kaplan, L. Kozma, O. Zamir, U. Zwick, Selection from Heaps, Row-Sorted Matrices, and $X+Y$ Using Soft Heaps, SOSA 2019.
\bibitem{barvinok} A. I. Barvinok, A Polynomial Time Algorithm for Counting Integral Points in Polyhedra When the Dimension is Fixed, Mathematics of Operations Research, 19(4), 1994.
\end{thebibliography}
\end{document}
